

\chapter{Conclusions}

In this project, we successfully improved the previous work of 
\cite{10.2312:ceig.20241139}, adding a mechanical model to simulate 
the physical properties of terrains and better reproduce erosion 
processes. Moreover, we also made improvements in the voronoi tessellation and 
added a low-resolution layer to speed up the erosion simulation. 
All added improvements with respect to the previous project contributed 
to the quality and realism of the final eroded terrain, which 
has the potential to be used to produce realistic landscapes in virtual 
scenes. 

\vspace{10pt}

Despite all of this, there is some work that can still be done to improve 
this project. First, the goal of this thesis was just to improve 
the previous method and add the physical considerations. In this 
project, we did not have the intention to provide tools for authoring 
terrains using our technique, and we just provided some tools for 
experimentation and debugging. If we want this method to be available for 
digital artists, we would need to port our algorithm to another application 
better suited for these tasks. Furthermore, the current lighting 
method of our application is not suited to produce high-quality 
visuals, so a port to another application that does implement 
better lighting will also improve this aspect.

\vspace{10pt}

Second, we have researched a little about improvements on the solver 
for systems of equations. There are some steps that can be done to improve 
that part, like using a GPU solver or recomputing just the relevant 
parts of the system when a portion of the terrain changes, avoiding a 
full computation each time. We did not have the time to test these improvements, 
so we leave them for future work. Additionally, even if the system of 
equations is the heaviest part of the project, it is not the only one 
that can be optimized; multiple parts of the code can be improved and, 
especially, parallelized in order to increase the algorithm efficiency.

\vspace{10pt}

The last potential improvement of this project has to do with parameter tuning. 
This thesis proposed multiple improvements, so we had to test a wide 
variety of parameters and methods to have a good understanding of how 
our techniques work and how they are compared against other methods. A 
possible extension of our work could consist in testing and comparing 
more materials and better tuning the parameters that control the link 
damage.  

\vspace{10pt}

Finally, and to end this project, we would like to thank both 
Oscar Argudo for tutoring this thesis and Diego Mateos for 
providing the base method this project is based on. Without their 
contributions this project would have not been possible.